\documentclass[conference]{IEEEtran}
\IEEEoverridecommandlockouts
\usepackage{cite}
\usepackage{amsmath,amssymb,amsfonts}
\usepackage{algorithmic}
\usepackage{graphicx}
\usepackage{textcomp}
\usepackage{xcolor}
\usepackage{booktabs}
\usepackage{url}

\begin{document}

\title{SuperBox: An Open Marketplace for Secure Discovery, Verification,
and Sandboxed Execution of Model Context Protocol Servers}

\author{
  \IEEEauthorblockN{Areeb Ahmed Shariff}
  \IEEEauthorblockA{%
    \textit{Dept. of Computer Science \& Business Systems}\\
    \textit{Dayananda Sagar College of Engineering}\\
    Bengaluru, India\\
    areebahmed0709@gmail.com}
  \and
  \IEEEauthorblockN{Amartya Anand}
  \IEEEauthorblockA{%
    \textit{Dept. of Computer Science \& Business Systems}\\
    \textit{Dayananda Sagar College of Engineering}\\
    Bengaluru, India\\
    amartyaanand07@gmail.com}
  \and
  \IEEEauthorblockN{Arush Verma}
  \IEEEauthorblockA{%
    \textit{Dept. of Computer Science \& Business Systems}\\
    \textit{Dayananda Sagar College of Engineering}\\
    Bengaluru, India\\
    arush3218@gmail.com}
  \and
  \IEEEauthorblockN{Devansh Aryan}
  \IEEEauthorblockA{%
    \textit{Dept. of Computer Science \& Business Systems}\\
    \textit{Dayananda Sagar College of Engineering}\\
    Bengaluru, India\\
    devansh123aryan@gmail.com}
}

\maketitle

\begin{abstract}
The rapid adoption of Model Context Protocol (MCP) servers has allowed
AI applications to interact with external tools, databases, and data
sources, enabling context-aware AI assistants and autonomous agents. However, the MCP ecosystem faces major
challenges, including the lack of a central registry, inconsistent
security checks, and no standardized setup. Servers are currently discovered manually through scattered GitHub
repositories and run directly on developer machines, with no mechanism
to assess reliability, verify security compliance, or manage deployments
across cloud and agent environments. These gaps constrain enterprise
adoption and expose users to real security risks when integrating
third-party MCP servers.

SuperBox addresses these gaps as an open-source marketplace and
serverless execution platform for MCP servers. It enforces an automated
five-stage security pipeline comprising SonarCloud for static analysis,
Snyk for dependency vulnerability detection, GitGuardian for secret
exposure prevention, and Bandit for Python-specific security linting.
SuperBox pairs a high-performance Go REST API with a developer-facing
CLI for publishing, discovery, and client configuration. AWS Lambda and
WebSocket API Gateway provide fully isolated, per-session execution
environments with enforced memory limits, timeouts, and controlled
network access, supporting real-time JSON-RPC communication.

By combining serverless execution, automated security validation,
S3-backed discovery, and auto-generated client configuration, SuperBox
reduces the operational complexity of deploying MCP servers while
enforcing consistent security standards. The platform is validated through security pipeline evaluation,
execution performance benchmarks, and end-to-end publish-to-execution
cycle testing across representative MCP server configurations. SuperBox is released as an open-source
platform, providing the MCP ecosystem with a validated, production-ready
foundation for secure server discovery, verification, and execution.
\end{abstract}

\begin{IEEEkeywords}
Model Context Protocol, Serverless Computing, AWS Lambda, WebSocket API,
Security Automation, AI Agent Infrastructure, Cloud-Native Architecture,
JSON-RPC
\end{IEEEkeywords}

%-----------------------------------------------------------------------
\section{Introduction}
\label{sec:intro}
%-----------------------------------------------------------------------

The Model Context Protocol (MCP), introduced by Anthropic in late 2024,
defines a standardized bidirectional communication interface between
large language models and external tools, APIs, and data sources
\cite{anthropic2024}. Within
its first year of release, the protocol surpassed eight million weekly
SDK downloads, signaling rapid and widespread adoption across AI
development teams \cite{hou2025}. MCP enables AI applications to move beyond static
knowledge, allowing them to query live databases, invoke APIs, read file
systems, and execute code through a unified protocol layer. This shift
has made MCP servers a foundational component in the emerging ecosystem
of autonomous AI agents and context-aware assistants.

Despite this growth, the MCP ecosystem lacks the infrastructure required
for reliable, secure, and scalable server deployment. Developers
currently discover MCP servers through scattered GitHub repositories,
personal recommendations, and informal aggregator sites, with no central
registry enforcing any form of quality or security standard. Most MCP
servers are consumed directly from source with no scanning, no
validation, and no isolation, meaning the code runs directly on the
developer's machine with access to local resources. Empirical studies of
the ecosystem have found that more than half of listed servers are non-functional or
low-value \cite{guo2025}, that dependency monocultures create
systemic risks, and that malicious servers have been uploaded to
widely-used aggregation platforms without triggering any audit
mechanism \cite{song2025}. This creates a critical trust
gap that constrains enterprise adoption and exposes users to real
security risks including data exfiltration, credential theft, and tool
poisoning.

This paper presents SuperBox, an open-source marketplace and serverless
execution platform for MCP servers. Taking inspiration from the
registry-and-validation model established by Docker Hub for container
images, SuperBox adapts that paradigm to the MCP ecosystem and addresses
the above challenges through three integrated contributions. First, a centralized S3-backed
registry with enforced metadata standards provides a single, trusted
source for MCP server discovery. Second, an automated five-stage
security pipeline covering static analysis, dependency scanning, secret
detection, and tool discovery ensures that only validated servers enter
the registry. Third, a serverless execution layer built on AWS Lambda
and WebSocket API Gateway provides completely isolated, on-demand
sandboxed execution environments for every MCP server invocation.
Together, these components eliminate the manual, fragmented, and
insecure patterns that currently characterize MCP server integration.

The remainder of this paper is organized as follows:
Section~\ref{sec:related} reviews related work on MCP ecosystem security
and deployment challenges. Section~\ref{sec:design} describes the system
design and architecture. Section~\ref{sec:impl} covers the
implementation details. Section~\ref{sec:results} presents results and
evaluation. Section~\ref{sec:conclusion} concludes with future work
directions.

%-----------------------------------------------------------------------
\section{Related Work}
\label{sec:related}
%-----------------------------------------------------------------------

The security and reliability challenges of the MCP ecosystem have
attracted significant research attention since the protocol's release. A
comprehensive lifecycle analysis of MCP servers structures the problem
across four phases (creation, deployment, operation, and maintenance)
and identifies sixteen distinct threat scenarios across four attacker
categories including malicious developers, external attackers, and
security flaws inherent to the protocol itself \cite{jamshidi2025}. This
foundational taxonomy establishes that MCP's open and extensible design,
while powerful, introduces trust challenges at every phase of the server
lifecycle that cannot be addressed through runtime defenses alone.

Large-scale empirical measurement of the ecosystem provides concrete
evidence of the fragmentation problem. A systematic study aggregating
over 8,400 valid MCP projects from six major markets found that more
than half of listed projects are invalid or low-value, that servers
exhibit structural risks including dependency monocultures and uneven
maintenance, and that no centralized mechanism exists for enforcing
quality or trust at the point of publication \cite{guo2025}. A parallel
study of 1,899 open-source MCP servers using a hybrid static analysis
pipeline found eight distinct vulnerability classes unique to MCP:
7.2\% of servers contained general security vulnerabilities and 5.5\%
exhibited MCP-specific tool poisoning, neither detectable by traditional
static analysis tools \cite{hasan2025}. Pre-deployment scanning must
therefore be MCP-aware to surface threats that general-purpose tools
would miss.

On the attack surface, research has identified several novel threat
classes specific to MCP's architecture. Parasitic Toolchain Attacks exploit MCP's lack of context-tool isolation
by embedding malicious instructions in external data sources accessed
during legitimate LLM tasks; the attack was demonstrated across a census
of 12,230 tools from 1,360 real servers, enabling stealthy data
exfiltration \cite{zhao2025}. Preference Manipulation Attacks demonstrate
that customized MCP servers can be crafted to influence LLMs into
prioritizing them over competing servers, representing a commercially
motivated threat that exploits open registration models \cite{wang2025}.
Tool poisoning, shadowing, and rug pull attacks have been demonstrated
across multiple model architectures; layered defenses combining manifest
signing and semantic vetting show partial mitigation but at the cost of
increased latency \cite{jamshidi2025}.

An end-to-end empirical attack evaluation confirmed that malicious MCP
servers were successfully uploaded to three widely-used aggregation
platforms without triggering any audit mechanism, and that users
consistently failed to distinguish malicious servers from legitimate
ones \cite{song2025}.

Proposed solutions in the literature primarily target runtime and
protocol-level defenses. A Zero Trust registry architecture with
admin-controlled registration, dynamic trust scoring based on versioning
and known vulnerabilities, and just-in-time credential provisioning has
been proposed to counter tool squatting in multi-agent
systems \cite{narajala2025}. However, no existing platform integrates
pre-deployment security validation, centralized registry management,
and isolated execution into a single deployable system. SuperBox
addresses this gap by moving the security boundary upstream to the point
of publication, combining pre-deployment validation with sandboxed
execution to provide defense across the full server lifecycle.

%-----------------------------------------------------------------------
\section{System Design and Architecture}
\label{sec:design}
%-----------------------------------------------------------------------

\subsection{Overview}

SuperBox is composed of three loosely coupled layers: the publish and
validation layer operated through the developer CLI, the discovery and
management layer served by the Go REST API and S3 registry, and the
execution layer powered by AWS Lambda and WebSocket API Gateway. Each
layer operates independently, communicates through well-defined
interfaces, and can be scaled or replaced without affecting the others.
This separation of concerns ensures that the security pipeline, registry
management, and execution sandbox never share runtime state, minimizing
the blast radius of any failure or compromise in a single layer.

\subsection{Security Pipeline}

The pipeline is designed against a threat model in which the primary
adversary is a developer who intentionally or negligently publishes a
server containing malicious code, exposed secrets, vulnerable
dependencies, or tool poisoning payloads. Threats requiring runtime
network interception or hardware-level compromise of the execution
environment are considered out of scope.

When a developer initiates a publish workflow, the CLI sequentially
executes five validation stages against the target GitHub repository
before any data is written to the registry. SonarCloud performs static
code analysis, detecting code smells, bugs, security hotspots, and
cognitive complexity issues across the full repository. Snyk audits all
declared dependencies in \texttt{requirements.txt} and
\texttt{package.json} files against its vulnerability database,
identifying known CVEs and suggesting remediation paths. GitGuardian
scans all files and commit history for hardcoded secrets, API keys,
tokens, and credentials. Bandit performs Python-specific security
linting, flagging dangerous function calls, insecure configurations, and
injection-prone patterns. A regex-based tool discovery pass then
extracts the names and signatures of all MCP tools exposed by the server
from the source code.

Each stage produces a structured result object. If any stage reports a
critical finding, the pipeline halts immediately and returns a
consolidated report to the developer. Only repositories that clear all
five stages proceed to registration. Table~\ref{tab:pipeline} summarizes the coverage of each stage.

\begin{table}[htbp]
  \caption{Security Pipeline Stage Coverage}
  \label{tab:pipeline}
  \centering
  \begin{tabular}{llp{3.8cm}}
    \toprule
    \textbf{Stage} & \textbf{Tool} & \textbf{Scope} \\
    \midrule
    Static Analysis  & SonarCloud  & Code smells, bugs, hotspots, complexity \\
    Dependency Scan  & Snyk        & Known CVEs in declared dependencies \\
    Secret Detection & GitGuardian & Hardcoded keys, tokens, credentials \\
    Python Linting   & Bandit      & Dangerous calls, injection-prone patterns \\
    Tool Discovery   & Regex pass  & MCP tool names, signatures, decorators \\
    \bottomrule
  \end{tabular}
\end{table}

\subsection{Registry Design}

The registry is implemented as a flat-file store in an AWS S3
bucket, where each MCP server is represented as a single JSON object
stored at the bucket root under the server's name as the key. This
design eliminates the operational overhead of a relational database for
what is fundamentally a read-heavy workload, enables direct access from
both the Go API and the Lambda executor without an intermediate data
layer, and produces a registry format that is portable, auditable, and
inspectable without any tooling. Each registry record contains the
server name, repository reference, entrypoint file, runtime language,
list of discovered tools, tool count, full security report from all five
pipeline stages, pricing configuration, and creation and update
timestamps. Registry contents are exposed through authenticated REST endpoints
covering listing, search, retrieval, creation, update, and deletion.
Table~\ref{tab:registry} lists all fields in each registry record.

\begin{table}[htbp]
  \caption{S3 Registry JSON Record Structure}
  \label{tab:registry}
  \centering
  \begin{tabular}{lp{4.5cm}}
    \toprule
    \textbf{Field} & \textbf{Description} \\
    \midrule
    \texttt{name}           & Unique server identifier (registry key) \\
    \texttt{repo}           & GitHub repository URL \\
    \texttt{entrypoint}     & Entry file path within the repository \\
    \texttt{language}       & Runtime language (Python, Node.js) \\
    \texttt{tools}          & List of discovered MCP tool descriptors \\
    \texttt{tool\_count}    & Total number of exposed tools \\
    \texttt{security}       & Full five-stage pipeline report object \\
    \texttt{pricing}        & Tier and price configuration (free/paid) \\
    \texttt{created\_at}   & ISO\ 8601 timestamp of first publication \\
    \texttt{updated\_at}   & ISO\ 8601 timestamp of last update \\
    \bottomrule
  \end{tabular}
\end{table}

\subsection{Execution Layer}

The execution layer bridges two fundamentally different communication
models: the stdio-based JSON-RPC interface expected by MCP servers and
the WebSocket-based interface used by AI clients connecting to the cloud.
A lightweight local proxy process, started by the CLI during client
configuration, reads JSON-RPC messages from the AI client's stdin and
forwards them over a persistent WebSocket connection to the AWS WebSocket
API Gateway. The API Gateway routes the connection lifecycle events and
messages to the Lambda function handler.

The Lambda handler exposes three WebSocket route handlers managing the
session lifecycle: connection setup, per-message dispatch, and
disconnection cleanup. On the first inbound message, the handler
bootstraps a fresh execution environment, performs the MCP
initialization handshake automatically, and then forwards the original
request. All subsequent messages are forwarded directly. On disconnect,
the server subprocess is terminated and all temporary files are purged,
ensuring complete isolation between sessions.

\subsection{Architecture Summary}

The three layers interact through the following paths. The publish path
flows from the developer CLI through the five-stage security pipeline
into the S3 registry. The discovery path flows from web and API clients
through the Go REST API, which reads directly from S3. The execution
path flows from an AI client through a local proxy over WebSocket to the
API Gateway, then to Lambda, which bootstraps the server process and
streams responses back through the same connection.

%-----------------------------------------------------------------------
\section{Implementation}
\label{sec:impl}
%-----------------------------------------------------------------------

\subsection{CLI and Publish Workflow}

The CLI is implemented in Python 3.11 using the Click framework
and installed as a standalone package via pip. The CLI exposes nine
commands: \texttt{auth} for registration, login, logout, status, and
token refresh via Firebase; \texttt{init} to scaffold a new
\texttt{superbox.json} configuration file; \texttt{push} to execute the
full security pipeline and publish to the registry; \texttt{pull} to
fetch registry metadata and write AI client configuration; \texttt{run}
to start an interactive JSON-RPC test session against a deployed server;
\texttt{search} to query the registry by name or keyword;
\texttt{inspect} to display full server details and security report;
\texttt{test} to run a server locally from a repository URL without
registry; and \texttt{logs} to fetch and optionally stream CloudWatch
execution logs.

The \texttt{push} command is the core of the publish workflow. It reads
the local \texttt{superbox.json} configuration file to extract the
repository URL, entrypoint, and language. It verifies Firebase
authentication by checking the stored token against the Identity Toolkit
API. It then runs the five-stage security pipeline sequentially, cloning
the repository to a temporary directory for local scanning stages. Tool
discovery uses regex patterns to identify MCP tool decorator definitions
across Python source files and \texttt{package.json} definitions for
Node.js servers. On successful completion of all stages, the command
serializes the full scan results and server metadata into a JSON record
and uploads it to the S3 registry bucket using boto3.

Table~\ref{tab:cli} lists all nine commands.

\begin{table}[htbp]
  \caption{CLI Command Reference}
  \label{tab:cli}
  \centering
  \begin{tabular}{lp{5cm}}
    \toprule
    \textbf{Command} & \textbf{Description} \\
    \midrule
    \texttt{auth}    & Register, login, logout, token refresh via Firebase \\
    \texttt{init}    & Scaffold a new \texttt{superbox.json} config file \\
    \texttt{push}    & Run security pipeline and publish to registry \\
    \texttt{pull}    & Fetch server metadata and write client config \\
    \texttt{run}     & Start interactive JSON-RPC session against server \\
    \texttt{search}  & Query registry by name or keyword \\
    \texttt{inspect} & Display full server details and security report \\
    \texttt{test}    & Run server locally from repo URL without registry \\
    \texttt{logs}    & Fetch and stream CloudWatch execution logs \\
    \bottomrule
  \end{tabular}
\end{table}

\subsection{API Server}

The REST API is implemented in Go 1.26 using the Gin web framework. It
exposes versioned endpoints under \texttt{/api/v1} for server management
(\texttt{GET /servers}, \texttt{GET /servers/:name},
\texttt{POST /servers}, \texttt{PUT /servers/:name},
\texttt{DELETE /servers/:name}) and authentication
(\texttt{POST /auth/register}, \texttt{POST /auth/login},
\texttt{POST /auth/device}, \texttt{GET /auth/profile}). Firebase token
validation is applied as middleware on all authenticated routes,
verifying the JWT against the Firebase public key set. S3 operations are
delegated to a Python helper subprocess using boto3, which the Go server
invokes via \texttt{exec.Command} and communicates with through
stdin/stdout. CORS is enabled for all origins to support web client
access. The server loads all configuration from environment variables at
startup, failing gracefully with degraded health status if required
variables are missing. Table~\ref{tab:api} summarizes the full API
surface.

\begin{table}[htbp]
  \caption{API Endpoint Summary}
  \label{tab:api}
  \centering
  \begin{tabular}{llcp{3.5cm}}
    \toprule
    \textbf{Method} & \textbf{Endpoint} & \textbf{Auth} & \textbf{Purpose} \\
    \midrule
    GET    & /api/v1/servers          & No  & List all registered servers \\
    GET    & /api/v1/servers/:name    & No  & Get server details and security report \\
    POST   & /api/v1/servers          & Yes & Register a new server (CLI push) \\
    PUT    & /api/v1/servers/:name    & Yes & Update an existing server record \\
    DELETE & /api/v1/servers/:name    & Yes & Remove a server from the registry \\
    POST   & /api/v1/auth/register    & No  & Create a new user account \\
    POST   & /api/v1/auth/login       & No  & Authenticate and receive JWT \\
    POST   & /api/v1/auth/device      & No  & Initiate CLI device flow OAuth \\
    GET    & /api/v1/auth/profile     & Yes & Get authenticated user profile \\
    \bottomrule
  \end{tabular}
\end{table}

\subsection{Lambda Executor}

The execution handler is implemented in Python 3.11 and deployed
to AWS Lambda with 1,024\,MB memory and a 60-second timeout,
configurable up to 900 seconds \cite{jonas2019}. The handler maintains a module-level process reference
and connection parameter store to persist the running MCP server process
across multiple WebSocket messages within a single connection lifecycle.
The \texttt{\$connect} handler extracts and stores the query string
parameters including the server name and optional test mode flags. The
\texttt{\$default} handler checks whether the MCP process is running; if
not, it fetches the server record from S3, clones the GitHub repository
by downloading the ZIP archive and extracting it to \texttt{/tmp}, runs
\texttt{pip install -r requirements.txt} with the target directory set to
\texttt{/tmp/pip\_modules}, and starts the MCP server as a subprocess
with stdio pipes. It then performs the MCP initialize handshake and
\texttt{notifications/initialized} notification before processing the
original request. The \texttt{\$disconnect} handler kills the subprocess
and removes all temporary files. Responses are sent back to the connected
client using the \texttt{apigatewaymanagementapi} boto3 client with the
connection's domain name and stage reconstructed from the event context.

\subsection{Frontend}

The web marketplace is implemented in Next.js~16 with React~19,
TypeScript~5, and Tailwind CSS~4. It provides server listing and search
with grid and list views, detailed server pages showing tool definitions,
security report summaries, repository metadata, and pricing.
The interface also provides Firebase-authenticated user profiles, a
guided publish workflow modal, a Razorpay-integrated paywall for premium
servers, and an interactive playground for live server testing. The
frontend communicates with the Go API for all data operations and with
Firebase directly for authentication state management.

\subsection{Infrastructure}

All AWS infrastructure is provisioned and managed using Terraform with a
modular structure. The S3 module creates the registry bucket with
appropriate access policies. The IAM module creates the Lambda execution
role with \texttt{s3:GetObject}, \texttt{logs:CreateLogGroup},
\texttt{logs:CreateLogStream}, and \texttt{logs:PutLogEvents}
permissions. The Lambda module packages and deploys the handler with
environment variables and configures a Function URL for direct
invocation. The WebSocket module creates the API Gateway WebSocket API,
defines the three route integrations, creates a Lambda integration with
the handler, and deploys to a named stage. Outputs from each module are
passed as inputs to dependent modules, and the final WebSocket URL and
S3 bucket name are output for use in the application environment
configuration.

%-----------------------------------------------------------------------
\section{Results and Evaluation}
\label{sec:results}
%-----------------------------------------------------------------------

\subsection{Security Pipeline Evaluation}

The five-stage security pipeline was evaluated against a set of ten MCP
server repositories seeded with known vulnerability classes: hardcoded
API keys and database credentials, dependencies with known CVEs at
varying severity levels, Python-specific dangerous patterns such as
\texttt{eval()} and \texttt{subprocess} with shell injection risks, and
tool descriptions containing prompt injection payloads. Each
vulnerability class was represented by at least one seeded instance per
repository. The pipeline correctly identified and flagged every seeded
finding. SonarCloud detected all code quality and injection-risk
patterns. Snyk flagged all CVE-matched dependencies. GitGuardian
identified all hardcoded credential patterns. Bandit caught all
Python-specific dangerous patterns. No false-positive rejections were
observed on clean repositories.

The tool discovery stage successfully extracted tool definitions from all
test repositories, including servers using both decorator-based and
registration-based MCP tool definition patterns.

\subsection{Execution Performance}

Lambda cold start performance was measured as the total time from
WebSocket \texttt{\$connect} event to the first valid JSON-RPC response
from the bootstrapped MCP server, across server repositories of varying
size and dependency complexity. Table~\ref{tab:coldstart} summarizes the results.

\begin{table}[htbp]
  \caption{Lambda Cold Start Latency by Repository Profile}
  \label{tab:coldstart}
  \centering
  \begin{tabular}{lcc}
    \toprule
    \textbf{Repository Profile} & \textbf{Deps} & \textbf{Observed Latency} \\
    \midrule
    Minimal Python (no external deps) & 0     & $\sim$8\,s   \\
    Light Python (stdlib + 1--3 pkgs) & 1--3  & $\sim$25\,s  \\
    Medium Python (several packages)   & 5--15 & $\sim$55\,s  \\
    Heavy Python (large frameworks)   & 15+   & $\sim$120\,s \\
    \bottomrule
  \end{tabular}
\end{table}

These measurements confirm that cold start latency is dominated by
dependency installation time. Subsequent messages within the same
WebSocket session are forwarded with sub-100\,ms overhead, as the MCP
process remains alive for the connection duration. The Lambda container
reuse behavior reduces cold starts for high-frequency servers.

\subsection{Registry and API Performance}

The S3-backed flat-file registry demonstrated consistent read
performance for list and individual retrieval operations. The Go API
returned full server lists with security reports in under 200\,ms for
registries of up to 500 servers in testing. A DynamoDB-backed
equivalent with the same metadata schema returned responses in a
comparable range but introduced additional connection management
overhead and per-read capacity unit cost. The flat-file design
eliminates query planning overhead and enables direct S3
\texttt{GetObject} calls for individual server lookups, producing
predictable latency regardless of registry size.

\subsection{End-to-End Validation}

The complete publish-to-execution cycle was validated across ten MCP
server repositories spanning Python REST API wrappers, file system tools,
database query tools, and web scraping servers. All servers that cleared
the five-stage pipeline were confirmed to be discoverable through the API
immediately after publication, with tool lists returned by the Lambda
executor at runtime matching the tool definitions extracted during
publish-time tool discovery. Client configuration files generated by the
CLI for all supported AI clients (Claude Desktop, Cursor, VSCode,
Windsurf, and ChatGPT) were confirmed to correctly route JSON-RPC traffic through the
local proxy to the Lambda execution layer.

%-----------------------------------------------------------------------
\section{Conclusion and Future Work}
\label{sec:conclusion}
%-----------------------------------------------------------------------

SuperBox addresses a critical infrastructure gap in the MCP ecosystem by
providing an integrated platform for centralized discovery, automated
security validation, and sandboxed execution of MCP servers. The
platform demonstrates that pre-deployment security validation at the
registry level is both technically feasible and practically necessary.
A serverless execution model built on AWS Lambda and WebSocket API
Gateway provides the session isolation and scalability required for safe
MCP server integration in production environments.

The primary contribution of this work is the shift of the security
boundary upstream to the point of publication rather than relying on
runtime defenses after deployment. By enforcing the five-stage pipeline as a precondition for registry
entry, the trust boundary is established at publication time,
complementing the runtime defenses proposed in related
work \cite{jamshidi2025,narajala2025}. The S3-backed flat-file registry design provides a portable,
auditable, and operationally simple foundation for the marketplace that
scales without a dedicated database tier.

The current implementation carries several limitations that bound the
scope of these results. Execution support is restricted to Python
servers; Node.js and other runtimes are not yet handled, excluding a
significant portion of the MCP ecosystem. Cold start latency for
dependency-heavy servers can exceed two minutes, making the platform
unsuitable for latency-sensitive workloads without warm container
strategies. The tool discovery stage relies on regex pattern matching,
which may miss dynamically registered tools or tools defined outside
standard decorator patterns. The security pipeline gates publication
but does not continuously monitor deployed servers for newly disclosed
CVEs in their dependencies.

Future work will target four directions: extending Lambda execution to
Node.js runtimes to cover the JavaScript portion of the MCP ecosystem;
automated repository webhook re-scanning to keep security reports current
without manual re-publication; MCP server versioning with rollback
capability to address the rug pull threat class identified in the
literature \cite{jamshidi2025}; and community-contributed ratings and download metrics to
add social trust signals alongside the automated pipeline results.

%-----------------------------------------------------------------------
\section*{Acknowledgment}
%-----------------------------------------------------------------------

The authors thank Ms.~Sanjana Shankar, Assistant Professor, Department
of Computer Science \& Business Systems, Dayananda Sagar College of
Engineering, Bengaluru, for her guidance throughout this work.

\begin{thebibliography}{00}

\bibitem{anthropic2024}
  Anthropic,
  ``Model Context Protocol Specification,''
  Nov. 2024. [Online]. Available: \url{https://modelcontextprotocol.io/specification}

\bibitem{jonas2019}
  E. Jonas, J. Schleier-Smith, V. Sreekanti, C.-C. Tsai, A. Khandelwal,
  Q. Pu, V. Shankar, J. Carreira, K. Krauth, N. Yadwadkar,
  J. E. Gonzalez, R. A. Popa, I. Stoica, and D. A. Patterson,
  ``Cloud Programming Simplified: A Berkeley View on Serverless
  Computing,''
  \textit{Tech. Rep.} UCB/EECS-2019-3, EECS Dept., Univ. of California,
  Berkeley, Feb. 2019. [Online]. Available:
  \url{https://www2.eecs.berkeley.edu/Pubs/TechRpts/2019/EECS-2019-3.html}

\bibitem{jamshidi2025}
  S. Jamshidi, K. W. Nafi, A. M. Dakhel, N. Shahabi, F. Khomh,
  and N. Ezzati-Jivan,
  ``Securing the Model Context Protocol: Defending LLMs Against Tool
  Poisoning and Adversarial Attacks,''
  \textit{arXiv preprint}, arXiv:2512.06556 [cs.CR], Dec. 2025.
  [Online]. Available: \url{https://arxiv.org/abs/2512.06556}

\bibitem{guo2025}
  H. Guo, Y. Hao, Y. Zhang, M. Xu, P. Lv, J. Chen \textit{et al.},
  ``A Measurement Study of Model Context Protocol Ecosystem,''
  \textit{arXiv preprint}, arXiv:2509.25292 [cs.CY], Sep. 2025.
  [Online]. Available: \url{https://arxiv.org/abs/2509.25292}

\bibitem{zhao2025}
  S. Zhao, Q. Hou, Z. Zhan, Y. Wang, Y. Xie, Y. Guo \textit{et al.},
  ``Mind Your Server: A Systematic Study of Parasitic Toolchain Attacks on
  the MCP Ecosystem,''
  \textit{arXiv preprint}, arXiv:2509.06572 [cs.CR], Sep. 2025.
  [Online]. Available: \url{https://arxiv.org/abs/2509.06572}

\bibitem{hasan2025}
  M. M. Hasan, H. Li, E. Fallahzadeh, G. K. Rajbahadur, B. Adams,
  and A. E. Hassan,
  ``Model Context Protocol (MCP) at First Glance: Studying the Security
  and Maintainability of MCP Servers,''
  \textit{arXiv preprint}, arXiv:2506.13538 [cs.SE], Jun. 2025.
  [Online]. Available: \url{https://arxiv.org/abs/2506.13538}

\bibitem{song2025}
  H. Song, Y. Shen, W. Luo, L. Guo, T. Chen, J. Wang \textit{et al.},
  ``Beyond the Protocol: Unveiling Attack Vectors in the Model Context
  Protocol (MCP) Ecosystem,''
  \textit{arXiv preprint}, arXiv:2506.02040 [cs.CR], Jun. 2025.
  [Online]. Available: \url{https://arxiv.org/abs/2506.02040}

\bibitem{wang2025}
  Z. Wang, R. Zhang, Y. Liu, W. Fan, W. Jiang, Q. Zhao \textit{et al.},
  ``MPMA: Preference Manipulation Attack Against Model Context Protocol,''
  \textit{arXiv preprint}, arXiv:2505.11154 [cs.CR], May 2025.
  [Online]. Available: \url{https://arxiv.org/abs/2505.11154}

\bibitem{narajala2025}
  V. S. Narajala, K. Huang, and I. Habler,
  ``Securing GenAI Multi-Agent Systems Against Tool Squatting: A Zero
  Trust Registry-Based Approach,''
  \textit{arXiv preprint}, arXiv:2504.19951 [cs.CR], Apr. 2025.
  [Online]. Available: \url{https://arxiv.org/abs/2504.19951}

\bibitem{hou2025}
  X. Hou, Y. Zhao, S. Wang, and H. Wang,
  ``Model Context Protocol (MCP): Landscape, Security Threats, and Future
  Research Directions,''
  \textit{arXiv preprint}, arXiv:2503.23278 [cs.CR], Mar. 2025.
  [Online]. Available: \url{https://arxiv.org/abs/2503.23278}

\end{thebibliography}

\end{document}
